\documentclass[11pt]{article}
\usepackage[T1]{fontenc}
\usepackage{lmodern}
\usepackage[margin=1in]{geometry}
\usepackage{amsmath,amssymb,amsthm,mathtools}
\usepackage{microtype}
\usepackage[colorlinks=true,linkcolor=blue,urlcolor=blue,citecolor=blue]{hyperref}
\usepackage{enumitem}
\setlist{nosep}
\newtheorem{theorem}{Theorem}
\newtheorem{lemma}[theorem]{Lemma}
\newtheorem{proposition}[theorem]{Proposition}
\newtheorem{remark}[theorem]{Remark}
\DeclareMathOperator{\dist}{dist}
\newcommand{\RR}{\mathbb R}
\newcommand{\CC}{\mathbb C}
\newcommand{\HH}{H^2(\RR^2;\CC)}
\newcommand{\nm}[2]{\lVert #1\rVert_{#2}}
\newcommand{\GG}{\mathcal G}
\title{Orbital Stability of the Complete Minimizing Set\\for a Fourth-Order Cubic--Quintic Schr\"odinger Equation}
\author{Richard J. Reyes}
\date{October 26, 2026}
\begin{document}
\maketitle
\begin{abstract}
For the Hamiltonian flow
\(i\partial_tu=(\Delta+k^2)^2u-\beta |u|^2u+\gamma |u|^4u\)
on \(\RR^2\), with \(\beta,\gamma,k>0\), we give a complete
variational proof that the full mass-constrained minimizing set is orbitally
stable in \(H^2\) for \(0<m\le m_B\), where
\(m_B=3\beta/[\gamma\sqrt{k^4+3\beta^2/(32\gamma)}]\).
The argument proves negativity of the infimum, strict binding, compactness
modulo translations, and global well-posedness, then applies conservation
of mass and energy. Perturbations are not required to preserve mass exactly.
This concerns the entire minimizing set, without uniqueness, attraction,
or spectral nondegeneracy.
\end{abstract}

\section{Statement and notation}
Write \(X=\HH\), \(A=\Delta+k^2\), and
\begin{align*}
M(u)&=\nm{u}{2}^2,&D(u)&=\nm{\Delta u}{2}^2,&K(u)&=\nm{\nabla u}{2}^2,\\
T(u)&=\nm{Au}{2}^2=D(u)-2k^2K(u)+k^4M(u),&
Q(u)&=\nm{u}{4}^4,&S(u)&=\nm{u}{6}^6.
\end{align*}
Set
\[
\begin{aligned}
H(u)&=T(u)-\frac{\beta}{2}Q(u)+\frac{\gamma}{3}S(u),\qquad
I(m)=\inf_{M(u)=m}H(u),\\
\GG_m&=\{u\in X:M(u)=m,\ H(u)=I(m)\}.
\end{aligned}
\]
Let
\[
\rho=k^4+\frac{3\beta^2}{32\gamma},\qquad
\sigma=\sqrt\rho,\qquad m_B=\frac{3\beta}{\gamma\sigma}.
\]
For a nonempty subset \(E\subset X\), define
\(\dist_X(v,E)=\inf_{g\in E}\nm{v-g}{H^2}\).

\begin{theorem}[Existence, compactness, and full-set orbital stability]\label{thm:main}
For each \(0<m\le m_B\):
\begin{enumerate}[label=(\roman*)]
\item \(-\infty<I(m)<0\), and \(I(m)<I(a)+I(m-a)\) for all \(0<a<m\);
\item \(\GG_m\ne\varnothing\), and every mass-\(m\) minimizing sequence
has a subsequence converging strongly in \(H^2\) after spatial translations
to an element of \(\GG_m\);
\item every datum in \(X\) generates a unique global solution
\(u\in C(\RR;X)\), conserving both \(M\) and \(H\);
\item for every \(\varepsilon>0\) there exists \(\delta>0\) such that
\[
\dist_X(u_0,\GG_m)<\delta\quad\Longrightarrow\quad
\sup_{t\in\RR}\dist_X(u(t),\GG_m)<\varepsilon.
\]
\end{enumerate}
The conclusion in (iv) permits initial data with \(M(u_0)\ne m\).
\end{theorem}

\section{Coercivity and strict negativity}
\begin{lemma}[Coercivity at fixed mass]\label{lem:coercivity}
With \(C_0=3\beta^2/(16\gamma)\), every \(u\in X\) satisfies
\[
H(u)\ge T(u)-C_0M(u),\qquad
\nm{\Delta u}{2}\le \sqrt{H(u)+C_0M(u)}+k^2\sqrt{M(u)}.
\]
On a set of bounded mass and bounded-above energy, the \(H^2\) norm is
uniformly bounded.
\end{lemma}
\begin{proof}
For \(r\ge0\), completion of the square gives
\[
-\frac\beta2r^2+\frac\gamma3r^3
=\frac\gamma3r\left(r-\frac{3\beta}{4\gamma}\right)^2
-\frac{3\beta^2}{16\gamma}r.
\]
Integrating with \(r=|u|^2\) gives the first inequality. Since
\(\Delta u=Au-k^2u\), the second follows from the triangle inequality.
Finally \(\nm{\nabla u}{2}^2\le\nm{u}{2}\nm{\Delta u}{2}\)
by Fourier Cauchy--Schwarz, so the \(L^2\) and Laplacian bounds control
\(\nm{u}{H^2}\).
\end{proof}

\begin{lemma}[Strict negativity at every positive mass]\label{lem:negative}
For every \(m>0\), \(-C_0m\le I(m)<0\).
\end{lemma}
\begin{proof}
The lower bound is Lemma~\ref{lem:coercivity}. For \(a,b,R>0\), take
\[
f_{a,b}(x,y)=(\pi ab)^{-1/2}
  \exp\left(-\frac{x^2}{2a^2}-\frac{y^2}{2b^2}\right),\qquad
u_R(x,y)=\frac{\sqrt m}{R}f_{a,b}(x/R,y/R)e^{ikx}.
\]
Then \(M(u_R)=m\). A direct Gaussian calculation, using
\(A(e^{ikx}f)=e^{ikx}(\Delta f+2ik\partial_xf)\), yields
\[
H(u_R)=\frac{c_2}{R^2}+\frac{c_4}{R^4},
\]
where
\begin{align*}
c_2&=\frac{2mk^2}{a^2}-\frac{\beta m^2}{4\pi ab},\\
c_4&=m\left(\frac{3}{4a^4}+\frac{1}{2a^2b^2}
+\frac{3}{4b^4}\right)+\frac{\gamma m^3}{9\pi^2a^2b^2}>0.
\end{align*}
Choosing \(a/b>8\pi k^2/(\beta m)\) makes \(c_2<0\). The choice
\(R^2=-2c_4/c_2\) then gives
\(H(u_R)=-c_2^2/(4c_4)<0\). Thus \(I(m)<0\).
\end{proof}

\section{Energy per mass and strict binding}
The dilation \(u_\lambda(x)=\lambda u(\lambda x)\) preserves mass
\emph{in two dimensions}. Its energy is
\begin{equation}\label{eq:dilation}
H(u_\lambda)=\lambda^4\left(D(u)+\frac\gamma3S(u)\right)
-\lambda^2\left(2k^2K(u)+\frac\beta2Q(u)\right)+k^4M(u).
\end{equation}

\begin{lemma}[Exact dilation adjustment]\label{lem:dilation}
Any minimizing sequence \((u_n)\) at any mass \(a>0\) can be replaced
by another minimizing sequence \((w_n)\) at the same mass such that
\begin{equation}\label{eq:V}
D(w_n)-k^2K(w_n)-\frac\beta4Q(w_n)
+\frac\gamma3S(w_n)=0
\end{equation}
for every \(n\).
\end{lemma}
\begin{proof}
Put \(P=D(u)+\gamma S(u)/3>0\) and
\(B=2k^2K(u)+\beta Q(u)/2>0\) for \(u\ne0\).
Equation~\eqref{eq:dilation} is \(Pz^2-Bz+k^4a\)
with \(z=\lambda^2>0\); its minimum occurs at
\(z=B/(2P)>0\). Set \(w_n=(u_n)_{\lambda_n}\) with this
minimizing scale. Then
\(I(a)\le H(w_n)\le H(u_n)\to I(a)\).
Differentiation of the exact polynomial at \(\lambda=1\), now applied
to each adjusted profile, yields~\eqref{eq:V}.
No differentiability of the map \(\lambda\mapsto u_\lambda\) in \(H^2\)
is assumed.
\end{proof}

\begin{lemma}[Uniform nonlinear ratio at an adjusted profile]\label{lem:ratio}
If \(w\in X\setminus\{0\}\) satisfies~\eqref{eq:V}, then
\[
D(w)\le\rho M(w),\qquad
\frac{S(w)}{Q(w)}\le\frac{\sigma M(w)}4.
\]
\end{lemma}
\begin{proof}
Writing \(r=|w|^2\), substitution of~\eqref{eq:V} into the
identity \(T=D-2k^2K+k^4M\) gives the exact equality
\begin{equation}\label{eq:completed}
D(w)+T(w)+\frac{2\gamma}{3}
\int_{\RR^2}r\left(r-\frac{3\beta}{8\gamma}\right)^2dx
=\left(k^4+\frac{3\beta^2}{32\gamma}\right)M(w)=\rho M(w).
\end{equation}
All three terms on the left are nonnegative, so \(D(w)\le\rho M(w)\).
With the Fourier convention
\(f(x)=(2\pi)^{-2}\int e^{ix\cdot\xi}\widehat f(\xi)\,d\xi\),
Cauchy--Schwarz and polar integration show, for any \(\eta>0\),
\[
\nm{w}{\infty}^2\le
\frac{D(w)+\eta M(w)}{8\sqrt\eta},\qquad
\frac{1}{(2\pi)^2}\int_{\RR^2}\frac{d\xi}{|\xi|^4+\eta}
=\frac1{8\sqrt\eta}.
\]
Here \(D(w)>0\) for nonzero \(w\in L^2\).
Taking \(\eta=D(w)/M(w)\) yields
\(\nm{w}{\infty}^2\le\sqrt{M(w)D(w)}/4
\le\sigma M(w)/4\).
Finally \(S(w)\le\nm{w}{\infty}^2Q(w)\), with \(Q(w)>0\).
\end{proof}

\begin{lemma}[Strictly decreasing specific energy]\label{lem:e}
Let \(e(a)=I(a)/a\). If \(0<a<b\) and \(a+b<2m_B\), then
\begin{equation}\label{eq:quantitative}
e(b)\le\left[1+\frac{b-a}{a}
\left(1-\frac{a+b}{2m_B}\right)\right]e(a)<e(a).
\end{equation}
Consequently \(e\) is strictly decreasing on \((0,m_B]\).
\end{lemma}
\begin{proof}
Put \(t=b/a>1\), \(h=1-(a+b)/(2m_B)>0\), and choose
an adjusted mass-\(a\) minimizing sequence \((w_n)\) from
Lemma~\ref{lem:dilation}. Lemma~\ref{lem:ratio} gives
\[
\frac\gamma3(t+1)S(w_n)
\le\frac{\gamma(t+1)\sigma a}{12}Q(w_n)
=\frac\beta2\frac{a+b}{2m_B}Q(w_n).
\]
Since \(M(\sqrt t\,w_n)=b\), the amplitude-scaling identity yields
\begin{align*}
\frac{I(b)}t&\le\frac{H(\sqrt t\,w_n)}t\\
&=H(w_n)-(t-1)\left(\frac\beta2Q(w_n)
-\frac\gamma3(t+1)S(w_n)\right)\\
&\le H(w_n)-(t-1)\frac{\beta h}{2}Q(w_n).
\end{align*}
Because \(T\ge0\) and \(S\ge0\), one has
\(Q(w_n)\ge-2H(w_n)/\beta\). Therefore
\[
\frac{I(b)}t\le[1+(t-1)h]H(w_n).
\]
Taking \(n\to\infty\) and dividing by \(a\) proves the first
inequality in~\eqref{eq:quantitative}. The second is strict because
\(e(a)=I(a)/a<0\) by Lemma~\ref{lem:negative} and
\((b-a)h/a>0\). If \(0<a<b\le m_B\), then
\(a+b<2m_B\), so the last assertion follows.
\end{proof}

\begin{proposition}[Strict binding]\label{prop:binding}
For every \(0<m\le m_B\) and every \(0<a<m\),
\[
I(m)<I(a)+I(m-a).
\]
\end{proposition}
\begin{proof}
By Lemma~\ref{lem:e}, \(e(a)>e(m)\) and
\(e(m-a)>e(m)\). Multiplying by the positive masses and adding,
\[
I(a)+I(m-a)=ae(a)+(m-a)e(m-a)>me(m)=I(m).
\]
\end{proof}

\section{Compactness modulo translations}
\begin{proposition}[Strong compactness of minimizing sequences]\label{prop:compact}
Fix \(0<m\le m_B\). If \(M(u_n)=m\) and \(H(u_n)\to I(m)\), then
there exist a subsequence, translations \(y_n\in\RR^2\), and
\(g\in\GG_m\) such that
\(u_n(\cdot+y_n)\to g\) strongly in \(H^2(\RR^2)\).
\end{proposition}
\begin{proof}
By Lemma~\ref{lem:coercivity}, \((u_n)\) is bounded in \(H^2\).
We first rule out \emph{vanishing}, namely
\[
\sup_{y\in\RR^2}\int_{B(y,1)}|u_n|^2\,dx\longrightarrow0.
\]
Cover \(\RR^2\) by unit squares \(Q_j\) (each contained in a unit
ball after translation). The local two-dimensional Gagliardo--Nirenberg
estimate on a unit square states
\[
\int_{Q_j}|v|^4\,dx\le C
\left(\int_{Q_j}|v|^2\,dx\right)
\left(\int_{Q_j}(|\nabla v|^2+|v|^2)\,dx\right).
\]
Summation over \(j\) would give
\(Q(u_n)\le C(\sup_j\nm{u_n}{L^2(Q_j)}^2)
\nm{u_n}{H^1}^2\to0\).
As \(H^2(\RR^2)\hookrightarrow L^\infty\), we also get
\(S(u_n)\le\nm{u_n}{\infty}^2Q(u_n)\to0\). This would imply
\(I(m)=\lim H(u_n)\ge0\), contrary to Lemma~\ref{lem:negative}.

Thus suitable translations make \(v_n=u_n(\cdot+y_n)\) retain a
uniformly positive \(L^2\) mass in a fixed ball. Passing to a subsequence,
\(v_n\rightharpoonup v\) weakly in \(H^2\), strongly in
\(L^p_{\rm loc}\) for all finite \(p\), and almost everywhere.
The local mass lower bound implies \(v\ne0\). Put \(a=M(v)\in(0,m]\)
and \(r_n=v_n-v\).
Weak convergence in \(H^2\) and the Hilbert-space identities give
\begin{equation}\label{eq:masssplit}
M(r_n)=m-a+o(1),\qquad
T(v_n)=T(v)+T(r_n)+o(1).
\end{equation}
The Brezis--Lieb lemma, applicable at powers \(4\) and \(6\) by boundedness
and almost-everywhere convergence, gives
\[
Q(v_n)=Q(v)+Q(r_n)+o(1),\qquad
S(v_n)=S(v)+S(r_n)+o(1).
\]
Hence
\begin{equation}\label{eq:energysplit}
H(v_n)=H(v)+H(r_n)+o(1).
\end{equation}
If \(a<m\), let \(b=m-a>0\) and
\(\widetilde r_n=\sqrt{b/M(r_n)}\,r_n\). The multiplying factor
converges to \(1\), and \((r_n)\) is uniformly bounded in \(H^2\).
Since \(H\) is a continuous polynomial functional with locally bounded
derivative on \(H^2\),
\(H(\widetilde r_n)-H(r_n)\to0\). Thus
\[
I(m)=\lim H(v_n)=H(v)+\lim H(r_n)
\ge I(a)+I(b),
\]
contradicting Proposition~\ref{prop:binding}.
Therefore \(a=m\).

Now \(M(r_n)\to0\). The uniform \(H^2\) bound implies
\(Q(r_n)+S(r_n)\to0\), since
\(Q(r_n)\le\nm{r_n}{\infty}^2M(r_n)\) and
\(S(r_n)\le\nm{r_n}{\infty}^4M(r_n)\).
Equation~\eqref{eq:energysplit} becomes
\(I(m)=H(v)+\lim T(r_n)\), while \(H(v)\ge I(m)\) and
\(T(r_n)\ge0\). It follows that \(H(v)=I(m)\) and \(T(r_n)\to0\).
Finally,
\[
\nm{\Delta r_n}{2}\le \nm{Ar_n}{2}+k^2\nm{r_n}{2}\to0,
\quad
\nm{\nabla r_n}{2}^2\le\nm{r_n}{2}\nm{\Delta r_n}{2}\to0.
\]
Thus \(r_n\to0\) strongly in \(H^2\), and \(v\in\GG_m\).
\end{proof}

\begin{remark}
The argument does not assume a radial, periodic, or real-valued class.
Translations are essential: \(\GG_m\) is not generally compact in the
original coordinates.
\end{remark}

\section{The global Hamiltonian flow}
\begin{proposition}[Global well-posedness and conservation]\label{prop:global}
For each \(u_0\in X\), the Cauchy problem
\begin{equation}\label{eq:flow}
i\partial_tu=A^2u-\beta |u|^2u+\gamma|u|^4u,
\qquad u(0)=u_0
\end{equation}
has a unique \(u\in C(\RR;X)\), depends continuously on data on
bounded time intervals, and satisfies
\(M(u(t))=M(u_0)\) and \(H(u(t))=H(u_0)\)
for all \(t\in\RR\).
\end{proposition}
\begin{proof}
The symbol of \(A^2\) is the real function
\((k^2-|\xi|^2)^2\). Therefore
\(U(t)=e^{-itA^2}\) is a strongly continuous unitary group on every
\(H^s\). In dimension two, \(H^2\hookrightarrow L^\infty\) and
\(H^2\) is a Banach algebra. Consequently the map
\(F(u)=-\beta|u|^2u+\gamma|u|^4u\) is locally Lipschitz
\(X\to X\). The Duhamel equation
\[
u(t)=U(t)u_0-i\int_0^t U(t-s)F(u(s))\,ds
\]
has, by contraction on \(C([-\tau,\tau];X)\) for small enough
\(\tau\), a unique maximal solution and the usual blow-up alternative:
if a finite positive or negative endpoint of the lifespan exists,
\(\nm{u(t)}{H^2}\) becomes unbounded there.

For completeness, conservation is justified at the energy regularity.
If \(u_0\in H^4\), the same contraction gives an \(H^4\) solution,
and \(u_t\in L^2\). Pairing~\eqref{eq:flow} with \(u\) and taking real
parts shows \(dM(u(t))/dt=0\), because \(A^2\) is self-adjoint and
both nonlinear pairings are real. The real Fr\'echet derivative of \(H\)
obeys
\[
DH(u)[h]=2\operatorname{Re}\int_{\RR^2}
\overline{\big(A^2u-\beta|u|^2u+\gamma|u|^4u\big)}\,h\,dx.
\]
Thus \(dH(u(t))/dt=2\operatorname{Re}\langle i u_t,u_t\rangle=0\).
(One may first use smoother data and pass to \(H^4\).)

For general \(u_0\in X\), choose \(u_{0,j}\in H^4\) converging in
\(H^2\). On any compact interval contained in the maximal \(H^2\)
lifespan, the local continuous-dependence estimates imply convergence
of the corresponding smooth solutions in \(C_tH^2_x\). Smooth solutions
continue in \(H^4\) throughout such an interval: the standard tame
product bound
\(\nm{F(u)}{H^4}\le C(\nm{u}{H^2})\nm{u}{H^4}\)
and Gronwall's inequality prevent an \(H^4\) blow-up while \(H^2\)
is bounded. Since \(M\) and \(H\) are continuous on \(H^2\), their
conservation passes to the limit.

Finally, Lemma~\ref{lem:coercivity} and the conservation laws give
\[
\nm{Au(t)}{2}^2\le H(u_0)+C_0M(u_0),\qquad
\nm{\Delta u(t)}{2}\le
\sqrt{H(u_0)+C_0M(u_0)}+k^2\sqrt{M(u_0)}.
\]
Together with fixed \(L^2\) mass, these yield a time-independent
\(H^2\) bound. The blow-up alternative excludes both finite-time
endpoints; hence the solution exists for all \(t\in\RR\).
\end{proof}

\section{Orbital stability of the full minimizing set}
\begin{proof}[Proof of Theorem~\ref{thm:main}(iv)]
Fix \(m\in(0,m_B]\). The first two claims of the theorem were proved
in Lemma~\ref{lem:negative}, Proposition~\ref{prop:binding}, and
Proposition~\ref{prop:compact}; the third is Proposition~\ref{prop:global}.
Moreover \(\GG_m\) is invariant under the global flow by conservation of
mass and energy, and is invariant under spatial translations by the
translation invariance of \(M\) and \(H\).

First note that \(\GG_m\) is uniformly bounded in \(H^2\): for any
\(g\in\GG_m\), Lemma~\ref{lem:coercivity} gives
\(T(g)\le I(m)+C_0m\), so
\(\nm{g}{H^2}\le C(m,\beta,\gamma,k)\).

Suppose stability fails. There are \(\varepsilon_0>0\) and, for
all sufficiently large \(n\), data \(u_{0,n}\in X\) for which
\(\dist_X(u_{0,n},\GG_m)<1/n\) while the supremum of the
distance along the flow is at least \(\varepsilon_0\).
For each \(n\), select \(t_n\in\RR\) with distance at least
\(\varepsilon_0/2\); this does not require the supremum to be attained. Then
\begin{equation}\label{eq:counter}
\dist_X(u_{0,n},\GG_m)\to0,\qquad
\dist_X(u_n(t_n),\GG_m)\ge\varepsilon_0/2,
\end{equation}
where \(u_n(t)\) is the global solution launched from \(u_{0,n}\).
Choose \(g_n\in\GG_m\) with
\(\nm{u_{0,n}-g_n}{H^2}\to0\). By the uniform bound on \(g_n\),
\((u_{0,n})\) is bounded in \(H^2\). Continuity of mass and energy,
uniform on bounded \(H^2\) sets for these polynomial functionals, yields
\begin{equation}\label{eq:near}
M(u_{0,n})\to m,\qquad H(u_{0,n})\to I(m).
\end{equation}

Set \(v_n=u_n(t_n)\). Conservation transfers~\eqref{eq:near} to
\(M(v_n)\to m\) and \(H(v_n)\to I(m)\).
Lemma~\ref{lem:coercivity} ensures \((v_n)\) is bounded in \(H^2\).
For large \(n\), define the \emph{mass-normalized} functions
\[
\alpha_n=\sqrt{\frac{m}{M(v_n)}},\qquad
w_n=\alpha_n v_n.
\]
Then \(\alpha_n\to1\), \(M(w_n)=m\), and
\[
\nm{w_n-v_n}{H^2}
=|\alpha_n-1|\nm{v_n}{H^2}\to0.
\]
The polynomial form of \(H\) and boundedness in \(H^2\) imply
\(H(w_n)-H(v_n)\to0\), so
\[
M(w_n)=m,\qquad H(w_n)\to I(m).
\]
By Proposition~\ref{prop:compact}, after passing to a subsequence there
exist \(y_n\in\RR^2\) and \(g\in\GG_m\) such that
\(w_n(\cdot+y_n)\to g\) strongly in \(H^2\).
Equivalently, \(\nm{w_n-g(\cdot-y_n)}{H^2}\to0\).
As \(g(\cdot-y_n)\in\GG_m\), we obtain
\[
\dist_X(v_n,\GG_m)\le
\nm{v_n-w_n}{H^2}
+\nm{w_n-g(\cdot-y_n)}{H^2}\longrightarrow0.
\]
This contradicts~\eqref{eq:counter}. Therefore orbital stability holds
for every \(\varepsilon>0\), including mass-changing perturbations.
\end{proof}

\section{Interpretation and antecedents}
The conclusion is \emph{Lyapunov stability of a set}. It does not imply
convergence as \(|t|\to\infty\), uniqueness of the minimizing profile,
stability of an arbitrary stationary solution outside \(\GG_m\), or
spectral nondegeneracy. The proof is specifically for two spatial
dimensions and does not classify four-spatial-dimensional lifetimes.

The final conservation-plus-compactness step follows the classical
Cazenave--Lions variational stability principle. Its application here
requires the preceding specific binding and compactness estimates for the
mixed cubic--quintic Hamiltonian. Closely related fourth-order single-power
mixed-dispersion problems were treated in the papers listed below.

\begin{thebibliography}{9}
\bibitem{CL82}
T. Cazenave and P.-L. Lions, ``Orbital stability of standing waves for some
nonlinear Schr\"odinger equations,'' \emph{Communications in Mathematical
Physics} \textbf{85} (1982), 549--561.
\href{https://doi.org/10.1007/BF01403504}{doi:10.1007/BF01403504}.
\bibitem{FJMM22}
A. J. Fern\'andez, L. Jeanjean, R. Mandel and M. Mari\c{s},
``Non-homogeneous Gagliardo--Nirenberg inequalities in
\(\mathbb R^N\) and application to a biharmonic non-linear Schr\"odinger equation,''
\emph{Journal of Differential Equations} \textbf{330} (2022), 1--65.
\href{https://doi.org/10.1016/j.jde.2022.04.037}{doi:10.1016/j.jde.2022.04.037}.
\bibitem{BCSN18}
D. Bonheure, J.-B. Casteras, E. Moreira dos Santos and R. Nascimento,
``Orbitally stable standing waves of a mixed dispersion nonlinear Schr\"odinger equation,''
\emph{SIAM Journal on Mathematical Analysis} \textbf{50} (2018), 5027--5071.
\href{https://doi.org/10.1137/17M1154138}{doi:10.1137/17M1154138}.
\bibitem{repository}
R. J. Reyes, ``Energy per mass,'' ``Strict subadditivity,'' ``Whole-space
compactness,'' and ``Orbital stability status,'' in the
\href{https://github.com/rickyjreyes/geometry_of_resonance/tree/main/docs/d4_binding}{geometry\_of\_resonance}
repository, \texttt{docs/d4\_binding/}, accessed October 26, 2026.
\end{thebibliography}
\end{document}